% 分类目录：ml
\documentclass[12pt, a4paper]{article}
\usepackage[margin=2.5cm]{geometry}
\usepackage{ctex}
\usepackage{amsmath, amssymb}
\usepackage{listings}
\usepackage{xcolor}
\usepackage{tabularx}
\usepackage{hyperref}

\lstset{
    language=Python,
    basicstyle=\ttfamily\small,
    keywordstyle=\color{blue},
    commentstyle=\color{green!50!black},
    stringstyle=\color{red},
    showstringspaces=false,
    breaklines=true,
    numbers=left,
    numberstyle=\tiny\color{gray},
    frame=single
}

\title{知识点：假设检验 (Hypothesis Testing)}
\author{}
\date{}

\begin{document}
\maketitle

\section{核心定义}
\textbf{中文定义}：假设检验是推断统计学中的一种决策过程，用于根据样本数据判断关于总体的某个声明（假设）是否成立。其基本思想是“小概率原理”：即在原假设成立的条件下，观测到目前样本数据的概率极低，从而怀疑原假设的真实性。

\textbf{英文定义}：Hypothesis testing is a formal procedure in inferential statistics used to determine whether there is enough evidence in a sample of data to infer that a certain condition is true for the entire population. It relies on the principle of rejecting a null hypothesis if the observed data is highly unlikely under that hypothesis.

\section{基本概念与术语}

\begin{itemize}
    \item \textbf{原假设 (Null Hypothesis, $H_0$)}：通常是希望通过检验推翻的陈述（如：新旧算法表现无显著差异）。
    \item \textbf{备择假设 (Alternative Hypothesis, $H_1$)}：原假设不成立时所接受的陈述。
    \item \textbf{显著性水平 ($\alpha$)}：拒绝原假设的“临界点”，常用的有 $0.05$ 或 $0.01$。它代表了犯第一类错误的概率。
    \item \textbf{P 值 (P-value)}：在 $H_0$ 成立的条件下，观测到当前统计量或更极端情况的概率。$P < \alpha$ 则拒绝 $H_0$。
    \item \textbf{第一类错误 (Type I Error)}：弃真错误，即 $H_0$ 为真却拒绝了它（概率为 $\alpha$）。
    \item \textbf{第二类错误 (Type II Error)}：取伪错误，即 $H_0$ 为假却接受了它（概率为 $\beta$）。
    \item \textbf{效能 (Power)}：$1 - \beta$，即当 $H_1$ 为真时，能正确拒绝 $H_0$ 的概率。
\end{itemize}

\section{基本流程}
\begin{enumerate}
    \item 建立原假设 $H_0$ 和备择假设 $H_1$。
    \item 选择显著性水平 $\alpha$ 和检验统计量（如 $Z, t, \chi^2$）。
    \item 计算检验统计量的观测值。
    \item 确定拒绝域或计算 P 值。
    \item 做出判断：若 $P \le \alpha$，拒绝 $H_0$；否则不拒绝 $H_0$。
\end{enumerate}

\section{常用检验方法}

\subsection{T 检验 (T-Test)}
适用于总体标准差未知且样本量较小（通常 $<30$）的情况。
\begin{itemize}
    \item \textbf{单样本 T 检验}：检验样本均值是否等于某个已知总体均值。
    \item \textbf{独立双样本 T 检验}：检验两个独立样本的均值是否存在显著差异。
    \item \textbf{配对样本 T 检验}：同一组对象在使用某种处理前后的差异对比。
\end{itemize}

\subsection{卡方检验 ($\chi^2$ Test)}
主要用于分类数据的分析。
\begin{itemize}
    \item \textbf{拟合优度检验}：观察频数是否符合某种理论分布。
    \item \textbf{独立性检验}：判断两个分类变量之间是否相关（如“吸烟”与“患肺癌”是否独立）。
\end{itemize}

\subsection{方差分析 (ANOVA)}
用于比较三个或更多组的均值是否存在显著差异（F 检验）。

\section{在机器学习中的应用}
\begin{itemize}
    \item \textbf{A/B 模拟测试}：对比实验组与对照组在转化率、点击率上的显著性差异。
    \item \textbf{特征选择}：通过卡方检验或 F 检验评估特征与标签的相关性。
    \item \textbf{模型比较}：对比两个模型在同一测试集上的表现差异（如 McNemar's Test 针对分类错误率）。
\end{itemize}

\section{关键代码片段 (Python)}
\begin{lstlisting}
from scipy import stats

# 1. 独立双样本 T 检验
group1 = [20, 22, 19, 21, 23]
group2 = [28, 30, 27, 29, 31]
t_stat, p_val = stats.ttest_ind(group1, group2)

# 2. 卡方独立性检验
observed = [[10, 20], [30, 40]]
chi2, p, dof, ex = stats.chi2_contingency(observed)

print(f"P-value: {p_val}")
if p_val < 0.05:
    print("拒绝原假设，差异显著")
\end{lstlisting}

\section{深度面试题}

\textbf{问：P 值的含义是什么？小 P 值一定代表效应很强吗？}\\
答：P 值是在原假设为真的前提下，得到当前结果或更极端结果的概率。P 值小只代表结果在统计学上是显着的，并不一定代表实际效应（Effect Size）很大。在大样本下，即使极微小的差异也能产生很小的 P 值。

\textbf{问：什么是第一类错误和第二类错误？为什么不能同时减小它们？}\\
答：第一类错误是“冤枉好人”，第二类错误是“放走坏人”。在固定样本量的情况下，减小 $\alpha$（更谨慎拒绝）必然会导致临界值向更极端移动，从而增大 $\beta$（更容易漏掉真实差异）。要同时减小两者，必须增加样本量。

\section{参考文献}
\begin{enumerate}
    \item Neyman, J., \& Pearson, E. S. (1933). \textit{On the Problem of the Most Efficient Tests of Statistical Hypotheses}.
    \item Fisher, R. A. (1925). \textit{Statistical Methods for Research Workers}.
    \item Wasserman, L. (2004). \textit{All of Statistics}.
\end{enumerate}

\end{document}
